\section{Conceptos en la modelación con incertidumbre}

Consideramos los siguientes conceptos que serán necesarios para desarrollar un modelo que actúe en un proceso estocástico.


\subsection{Probabilidad basica}


\begin{definition}
    Una variable aleatoria $X$ es una función de un espacio muestral $\Omega=\{s_1,s_2,s_3,\dots\}$ hacia los números reales. Con una funcion de probabilidad $\mathbb{P}$ en $\Omega$, sea $\mathcal{X}=\{x_1,\dots,x_m\}$ el rango de $X$. Definimos la funcion de probabilidad $\mathbb{P}_X$ en $\mathcal{X}$ observando que $X=x_i$ si y solo si el resultado de un experimento aleatorio es $s_j\in\Omega\;\text{t.q}\;X(s_j)=x_i$. Entonces:
    
    \begin{align}
        \mathbb{P}_X(X=x_i) = \mathbb{P}(\{s_j\in\Omega\,:\,X(s_j)=x_i\})
    \end{align}\autocite{:berger2002}.
\end{definition}

Para toda variable aleatoria $X$ le asociamos una función de distribución acumulada.

\begin{definition}
    Una funcion de distribución acumulada $F_X(x)$ es definida como:

    \begin{align}
        F_X(x)=\mathbb{P}_X(X\leq x)\;\forall\;x
    \end{align}\autocite{:berger2002}.
\end{definition}


Una variable aleatoria puede estar asociada a una función de probabilidad acumulada \textit{cdf} que puede ser o una función de densidad de probabilidad \textit{pdf} o una función de masa de probabilidad \textit{pmf}, esto dependiedo si esta en un espacio continuo o discreto respectivamente.

\begin{definition}
    Una \textit{pmf} para una variable aleatoria discreta es

    \begin{align}
        f_X(x)=\mathbb{P}_X(X=x)\,\forall\,x.
    \end{align}
\end{definition}

\begin{definition}
    Una \textit{pdf} de una variable aleatoria continua satisface

    \begin{align}
        F_X(x)=\int_{-\infty}^{x}f_X(t)\mathbf{dt}\,\forall\,x
    \end{align}\autocite{:berger2002}.
\end{definition}


Entre las distribuciones de probabilidad que podrian.